% GENERATED FILE. Edit canonical lessons, then run npm run generate:books.
% canonical-source-hash: fnv1a64:bcb221e216dbc574
% canonical-lessons: UR-C517-istej, UR-C517-tali, UR-C517-shaiqin, UR-C517-mukalama, UR-C517-lubhana, UR-R517-review-510-517

\chapter{On the Stage}
\label{ch:ur-on-the-stage}
\hlchaptermodality{517}

\begin{chapteropening}
\textbf{By the end of this chapter:} \emph{I can talk about a play and its audience.}

Complete the last lesson of chapter 517: I can talk about a play and its audience.
\end{chapteropening}

\section[\texorpdfstring{\ur{اسٹیج} (isṭej)}{اسٹیج (isṭej)}]{\ur{اسٹیج} (isṭej) — a stage (in a theatre)}

\label{lesson:UR-C517-istej}



\begin{quote}
Before the new one: say the Urdu for to engrave, then the Urdu for an idol.
\end{quote}

\subsection*{You'll want to know: \ur{اسٹیج}}
\textbf{\ur{اسٹیج}} — \emph{isṭej} — \textquotedblleft{}a stage (in a theatre)\textquotedblright{}.

\begin{quote}
\ur{اداکار} \ur{اسٹیج} \ur{پر} \ur{آیا۔}
\end{quote}

\begin{quote}
\emph{adākār isṭej par āyā.}
\end{quote}

\begin{quote}
The actor came onto the stage.
\end{quote}

Masculine.

\begin{cousinweb}
\textbf{\ur{اسٹیج}} (\emph{isṭej}) is borrowed from English \emph{stage}.
\end{cousinweb}

\begin{grammarlens}[title={What you've built}]
One more word for describing what you lived and felt.
\end{grammarlens}

\subsection*{Your turn}
\begin{itemize}
\raggedright
  \item \emph{Say it:} \emph{isṭej}
  \item \emph{Say it:} \emph{isṭej}, once more
  \item \emph{Say it:} \emph{adākār isṭej par āyā.} — The actor came onto the stage.
  \item \emph{Recall:} say the Urdu for to engrave, then the Urdu for an idol, then say \emph{isṭej} again
\end{itemize}

\begin{tcolorbox}[breakable,colback=teal!4,colframe=teal!35!black,title={Before you move on}]
What does \ur{اسٹیج} mean? (\textquotedblleft{}A stage (in a theatre)\textquotedblright{}.) Say it once more.
\end{tcolorbox}
\section[\texorpdfstring{\ur{تالی} (tālī)}{تالی (tālī)}]{\ur{تالی} (tālī) — a clap (of the hands)}

\label{lesson:UR-C517-tali}



\begin{quote}
Before the new one: say the Urdu for an idol, then the Urdu for a stage (in a theatre).
\end{quote}

\subsection*{You'll want to know: \ur{تالی}}
\textbf{\ur{تالی}} — \emph{tālī} — \textquotedblleft{}a clap (of the hands)\textquotedblright{}.

\begin{quote}
\ur{ناٹک} \ur{کے} \ur{بعد} \ur{سب} \ur{نے} \ur{تالیاں} \ur{بجائیں۔}
\end{quote}

\begin{quote}
\emph{nāṭak ke bād sab ne tāliyāṅ bajāīṅ.}
\end{quote}

\begin{quote}
After the play everyone clapped.
\end{quote}

Feminine.

\begin{cousinweb}
\textbf{\ur{تالی}} (\emph{tālī}) is inherited from Sanskrit \emph{tāla}, \textquotedblleft{}slapping the hands together; musical time\textquotedblright{} (Monier-Williams).
\end{cousinweb}

\begin{grammarlens}[title={What you've built}]
One more word for describing what you lived and felt.
\end{grammarlens}

\subsection*{Your turn}
\begin{itemize}
\raggedright
  \item \emph{Say it:} \emph{tālī}
  \item \emph{Say it:} \emph{tālī}, once more
  \item \emph{Say it:} \emph{nāṭak ke bād sab ne tāliyāṅ bajāīṅ.} — After the play everyone clapped.
  \item \emph{Recall:} say the Urdu for an idol, then the Urdu for a stage (in a theatre), then say \emph{tālī} again
\end{itemize}

\begin{tcolorbox}[breakable,colback=teal!4,colframe=teal!35!black,title={Before you move on}]
What does \ur{تالی} mean? (\textquotedblleft{}A clap (of the hands)\textquotedblright{}.) Say it once more.
\end{tcolorbox}
\section[\texorpdfstring{\ur{شائقین} (shāiqīn)}{شائقین (shāiqīn)}]{\ur{شائقین} (shāiqīn) — fans; keen spectators}

\label{lesson:UR-C517-shaiqin}



\begin{quote}
Before the new one: say the Urdu for a stage (in a theatre), then the Urdu for a clap.
\end{quote}

\subsection*{You'll want to know: \ur{شائقین}}
\textbf{\ur{شائقین}} — \emph{shāiqīn} — \textquotedblleft{}fans; keen spectators\textquotedblright{}.

\begin{quote}
\ur{شائقین} \ur{اداکاروں} \ur{کو} \ur{دیکھ} \ur{کر} \ur{ہنسے۔}
\end{quote}

\begin{quote}
\emph{shāiqīn adākāroṅ ko dekh kar hãse.}
\end{quote}

\begin{quote}
The audience laughed as they watched the actors.
\end{quote}

Masculine, used in the plural.

\begin{cousinweb}
\textbf{\ur{شائقین}} (\emph{shāiqīn}) is usually traced, through Persian, to Arabic \emph{shāʾiqīn}, the plural of \emph{shāʾiq}, \textquotedblleft{}eager, longing\textquotedblright{}: the people keen on a game or a show.
\end{cousinweb}

\begin{grammarlens}[title={What you've built}]
One more word for describing what you lived and felt.
\end{grammarlens}

\subsection*{Your turn}
\begin{itemize}
\raggedright
  \item \emph{Say it:} \emph{shāiqīn}
  \item \emph{Say it:} \emph{shāiqīn}, once more
  \item \emph{Say it:} \emph{shāiqīn adākāroṅ ko dekh kar hãse.} — The audience laughed as they watched the actors.
  \item \emph{Recall:} say the Urdu for a stage (in a theatre), then the Urdu for a clap, then say \emph{shāiqīn} again
\end{itemize}

\begin{tcolorbox}[breakable,colback=teal!4,colframe=teal!35!black,title={Before you move on}]
What does \ur{شائقین} mean? (\textquotedblleft{}Fans; keen spectators\textquotedblright{}.) Say it once more.
\end{tcolorbox}
\section[\texorpdfstring{\ur{مکالمہ} (mukālamā)}{مکالمہ (mukālamā)}]{\ur{مکالمہ} (mukālamā) — a line of dialogue; a dialogue}

\label{lesson:UR-C517-mukalama}



\begin{quote}
Before the new one: say the Urdu for a clap, then the Urdu for fans of a game or show.
\end{quote}

\subsection*{You'll want to know: \ur{مکالمہ}}
\textbf{\ur{مکالمہ}} — \emph{mukālamā} — \textquotedblleft{}a line of dialogue; a dialogue\textquotedblright{}.

\begin{quote}
\ur{اداکار} \ur{اپنے} \ur{مکالمے} \ur{بھول} \ur{گیا۔}
\end{quote}

\begin{quote}
\emph{adākār apne mukālame bhūl gayā.}
\end{quote}

\begin{quote}
The actor forgot his lines.
\end{quote}

Masculine.

\begin{cousinweb}
\textbf{\ur{مکالمہ}} (\emph{mukālamā}) is usually traced, through Persian, to Arabic \emph{mukālama}, \textquotedblleft{}talking together, a conversation\textquotedblright{}.
\end{cousinweb}

\begin{grammarlens}[title={What you've built}]
One more word for describing what you lived and felt.
\end{grammarlens}

\subsection*{Your turn}
\begin{itemize}
\raggedright
  \item \emph{Say it:} \emph{mukālamā}
  \item \emph{Say it:} \emph{mukālamā}, once more
  \item \emph{Say it:} \emph{adākār apne mukālame bhūl gayā.} — The actor forgot his lines.
  \item \emph{Recall:} say the Urdu for a clap, then the Urdu for fans of a game or show, then say \emph{mukālamā} again
\end{itemize}

\begin{tcolorbox}[breakable,colback=teal!4,colframe=teal!35!black,title={Before you move on}]
What does \ur{مکالمہ} mean? (\textquotedblleft{}A line of dialogue; a dialogue\textquotedblright{}.) Say it once more.
\end{tcolorbox}
\section[\texorpdfstring{\ur{لبھانا} (lubhānā)}{لبھانا (lubhānā)}]{\ur{لبھانا} (lubhānā) — to charm; to entice}

\label{lesson:UR-C517-lubhana}



\begin{quote}
Before the new one: say the Urdu for fans of a game or show, then the Urdu for a dialogue.
\end{quote}

\subsection*{You'll want to know: \ur{لبھانا}}
\textbf{\ur{لبھانا}} — \emph{lubhānā} — \textquotedblleft{}to charm; to entice\textquotedblright{}.

\begin{quote}
\ur{اس} \ur{کے} \ur{گیت} \ur{نے} \ur{سب} \ur{کو} \ur{لبھایا۔}
\end{quote}

\begin{quote}
\emph{us ke gīt ne sab ko lubhāyā.}
\end{quote}

\begin{quote}
His song charmed everyone.
\end{quote}

Past, with \ur{نے} (he or she …): \ur{اس} \ur{نے} \ur{سب} \ur{کو} \ur{لبھایا}, \emph{us ne sab ko lubhāyā}.

\begin{cousinweb}
\textbf{\ur{لبھانا}} (\emph{lubhānā}) is inherited: CDIAL lists the Hindi cousin \emph{lubhānā}, \textquotedblleft{}to excite desire\textquotedblright{}, under Sanskrit \emph{lubhyati}, \textquotedblleft{}desires\textquotedblright{} (CDIAL 11086).
\end{cousinweb}

\begin{grammarlens}[title={What you've built}]
One more word for describing what you lived and felt.
\end{grammarlens}

\subsection*{Your turn}
\begin{itemize}
\raggedright
  \item \emph{Say it:} \emph{lubhānā}
  \item \emph{Say it:} \emph{lubhānā}, once more
  \item \emph{Say it:} \emph{us ke gīt ne sab ko lubhāyā.} — His song charmed everyone.
  \item \emph{Recall:} say the Urdu for fans of a game or show, then the Urdu for a dialogue, then say \emph{lubhānā} again
\end{itemize}

\begin{tcolorbox}[breakable,colback=teal!4,colframe=teal!35!black,title={Before you move on}]
What does \ur{لبھانا} mean? (\textquotedblleft{}To charm; to entice\textquotedblright{}.) Say it once more.
\end{tcolorbox}
\section[(dialogue)]{Looking back: from Peel, Pound, Squeeze to On the Stage}

\label{lesson:UR-R517-review-510-517}



\begin{quote}
Say the Urdu for to peel and to pound.
\end{quote}

\subsection*{Your turn}
Say these aloud:

\begin{itemize}
\raggedright
  \item to peel, to pound, a mortar (for pounding), to have (something) ground, to squeeze out
  \item a recipe, a cloth spread for a meal, to have (food) cooked, to blow (on), to give (someone) a taste
  \item a garden, a flower bed, fertiliser, to water (plants), to uproot
  \item free time, a pastime, recreation, a picnic, an album
  \item wool, to thread (a needle), to embroider, a thimble, a fringe
  \item a picture, a painter (an artist), a sketch, a masterpiece, to show through
  \item calligraphy, carved or painted decoration, a figurine, to engrave, an idol
  \item a stage (in a theatre), a clap, fans of a game or show, a dialogue, to charm
\end{itemize}

\begin{tcolorbox}[breakable,colback=teal!4,colframe=teal!35!black,title={Before you move on}]
To peel? (\textbf{\ur{چھیلنا}}, \emph{chhīlnā}.) To pound? (\textbf{\ur{کوٹنا}}, \emph{kūṭnā}.) A mortar (for pounding)? (\textbf{\ur{اوکھلی}}, \emph{okhlī}.) To have (something) ground? (\textbf{\ur{پسوانا}}, \emph{pisvānā}.) To squeeze out? (\textbf{\ur{نچوڑنا}}, \emph{nichoṛnā}.) A recipe? (\textbf{\ur{ترکیب}}, \emph{tarkīb}.) A cloth spread for a meal? (\textbf{\ur{دسترخوان}}, \emph{dastarkhvān}.) To have (food) cooked? (\textbf{\ur{پکوانا}}, \emph{pakvānā}.) To blow (on)? (\textbf{\ur{پھونکنا}}, \emph{phūṅknā}.) To give (someone) a taste? (\textbf{\ur{چکھانا}}, \emph{chakhānā}.) A garden? (\textbf{\ur{باغ}}, \emph{bāgh}.) A flower bed? (\textbf{\ur{کیاری}}, \emph{kyārī}.) Fertiliser? (\textbf{\ur{کھاد}}, \emph{khād}.) To water (plants)? (\textbf{\ur{سینچنا}}, \emph{sīṅchnā}.) To uproot? (\textbf{\ur{اکھاڑنا}}, \emph{ukhāṛnā}.) Free time? (\textbf{\ur{فرصت}}, \emph{fursat}.) A pastime? (\textbf{\ur{مشغلہ}}, \emph{mashghalā}.) Recreation? (\textbf{\ur{تفریح}}, \emph{tafrīh}.) A picnic? (\textbf{\ur{پکنک}}, \emph{piknik}.) An album? (\textbf{\ur{البم}}, \emph{album}.) Wool? (\textbf{\ur{اون}}, \emph{ūn}.) To thread (a needle)? (\textbf{\ur{پرونا}}, \emph{pironā}.) To embroider? (\textbf{\ur{کاڑھنا}}, \emph{kāṛhnā}.) A thimble? (\textbf{\ur{انگشتانہ}}, \emph{angushtānā}.) A fringe? (\textbf{\ur{جھالر}}, \emph{jhālar}.) A picture? (\textbf{\ur{تصویر}}, \emph{tasvīr}.) A painter (an artist)? (\textbf{\ur{مصور}}, \emph{musavvir}.) A sketch? (\textbf{\ur{خاکہ}}, \emph{khākā}.) A masterpiece? (\textbf{\ur{شاہکار}}, \emph{shāhkār}.) To show through? (\textbf{\ur{جھلکنا}}, \emph{jhalaknā}.) Calligraphy? (\textbf{\ur{خطاطی}}, \emph{khattātī}.) Carved or painted decoration? (\textbf{\ur{نقاشی}}, \emph{naqqāshī}.) A figurine? (\textbf{\ur{مورتی}}, \emph{mūrtī}.) To engrave? (\textbf{\ur{کندہ} \ur{کرنا}}, \emph{kandā karnā}.) An idol? (\textbf{\ur{بت}}, \emph{but}.) A stage (in a theatre)? (\textbf{\ur{اسٹیج}}, \emph{isṭej}.) A clap? (\textbf{\ur{تالی}}, \emph{tālī}.) Fans of a game or show? (\textbf{\ur{شائقین}}, \emph{shāiqīn}.) A dialogue? (\textbf{\ur{مکالمہ}}, \emph{mukālamā}.) To charm? (\textbf{\ur{لبھانا}}, \emph{lubhānā}.)
\end{tcolorbox}
